\usetikzlibrary{fit, calc, arrows.meta, positioning, decorations.pathreplacing}

%%% ================================================================
%%% OPERATORS  (small inline operation nodes)
%%% ================================================================

\tikzset{
  op base/.style={draw, minimum size=0.35cm, inner sep=0pt},
  NOOP/.style={op base, draw=none, fill=none,},
  %% XOR: circle with a cross
  XOR/.style={op base, circle,
    append after command={
      [shorten >=\pgflinewidth, shorten <=\pgflinewidth]
      (\tikzlastnode.north) edge (\tikzlastnode.south)
      (\tikzlastnode.east)  edge (\tikzlastnode.west)}},
  %% AND: circle with x (boolean AND)
  AND/.style={op base, circle,
    append after command={
      [shorten >=\pgflinewidth, shorten <=\pgflinewidth]
      (\tikzlastnode.north west) edge (\tikzlastnode.south east)
      (\tikzlastnode.north east) edge (\tikzlastnode.south west)}},
  %% PLUS: square with + (use for any field/modular addition)
  PLUS/.style={op base, rectangle,
    append after command={
      [shorten >=\pgflinewidth, shorten <=\pgflinewidth]
      (\tikzlastnode.north) edge (\tikzlastnode.south)
      (\tikzlastnode.east)  edge (\tikzlastnode.west)}},
  %% MINUS: square with -
  MINUS/.style={op base, rectangle,
    append after command={
      [shorten >=\pgflinewidth, shorten <=\pgflinewidth]
      (\tikzlastnode.east) edge (\tikzlastnode.west)}},
  %% MUL: square with x
  MUL/.style={op base, rectangle,
    append after command={
      [shorten >=\pgflinewidth, shorten <=\pgflinewidth]
      (\tikzlastnode.north west) edge (\tikzlastnode.south east)
      (\tikzlastnode.north east) edge (\tikzlastnode.south west)}},
}


%%% ================================================================
%%% BOXES  (labeled function boxes, reused for S-box, F, pi, layers, ...)
%%% ================================================================

\tikzset{
  %% semantic fill colors -- apply to any box
  crypto role/.is choice,
  crypto role/default/.style = {fill=gray!10},
  crypto role/sbox/.style    = {fill=orange!15},
  crypto role/linear/.style  = {fill=blue!10},
  crypto role/key/.style     = {fill=green!15},
  crypto role/perm/.style    = {fill=purple!10},
  %% base box -- inherit this, don't use directly
  box base/.style={draw, inner sep=0.5pt, text height=2ex, align=center, crypto role=default},
  bbox/.style 2 args={box base, minimum width=#1, minimum height=#2},
  %% square box: S-box, single-wire function, constant
  sbox/.style={bbox={0.8cm}{0.8cm}},
  %% layer box: sizes itself to the bounding box of a given list of nodes
  %% usage: \node[layerbox={(a)(b)(c)}, right=of refnode] (name) {label};
  %layerbox/.style={box base, fit={#1}},
  layerbox/.style 2 args={box base, fit={#1}, label=center:{#2}},
  %% convenience aliases
  vbox/.style = {bbox={0.7cm}{#1}},
  hbox/.style = {bbox={#1}{0.7cm}},
  note/.style = {color=gray, font=\scriptsize},
}

%%% ================================================================
%%% STATE WIRES
%%% ================================================================

\tikzset{
  state/.style = {
    minimum size=0.5cm, inner sep=0pt, outer sep=0pt
  },
  wire/.style      = {-{Latex[length=5pt]},},
  wire dashed/.style = {wire, dashed},
}


%%% ================================================================
%%% BRACES
%%% ================================================================

%% \hbrace[left=, right=, raise=]{P1}{P2}{P3}
%%   horizontal brace at y(P3), from x(P1)-left to x(P2)+right
\pgfkeys{/hbrace/.is family, /hbrace,
  left/.initial=0pt, right/.initial=0pt, raise/.initial=4pt,
  color/.initial=gray, font/.initial=\scriptsize}
\newcommand{\hbrace}[4][]{%
  \pgfkeys{/hbrace, left=0pt, right=0pt, raise=4pt, color=gray, font=\scriptsize, #1}%
  \draw[decorate, decoration={brace, mirror},
        color=\pgfkeysvalueof{/hbrace/color}, font=\pgfkeysvalueof{/hbrace/font}]
    ($(#2 |- #4)+(-\pgfkeysvalueof{/hbrace/left},-\pgfkeysvalueof{/hbrace/raise})$) --
    ($(#3 |- #4)+(\pgfkeysvalueof{/hbrace/right},-\pgfkeysvalueof{/hbrace/raise})$)%
}

%% \vbrace[above=, below=, raise=]{P1}{P2}{P3}
%%   vertical brace to the right of x(P3)+raise, from y(P2)-below upward to y(P1)+above
%%   P1 = top node, P2 = bottom node, P3 = x reference (typically rightmost node)
%%   brace opens leftward toward content; node[midway, right=..] places label outside
\pgfkeys{/vbrace/.is family, /vbrace,
  above/.initial=0pt, below/.initial=0pt, raise/.initial=4pt,
  color/.initial=gray, font/.initial=\scriptsize}
\newcommand{\vbrace}[4][]{%
  \pgfkeys{/vbrace, above=0pt, below=0pt, raise=4pt, color=gray, font=\scriptsize, #1}%
  \draw[decorate, decoration={brace},
        color=\pgfkeysvalueof{/vbrace/color}, font=\pgfkeysvalueof{/vbrace/font}]
    ($(#3 -| #4)+(+\pgfkeysvalueof{/vbrace/raise},-\pgfkeysvalueof{/vbrace/below})$) --
    ($(#2 -| #4)+(+\pgfkeysvalueof{/vbrace/raise},+\pgfkeysvalueof{/vbrace/above})$)%
}


